\documentclass[10pt,letterpaper]{article}

\usepackage[T1]{fontenc}
\usepackage{lmodern}
\renewcommand{\familydefault}{\sfdefault}
\usepackage[letterpaper,top=0.68in,bottom=0.62in,left=0.78in,right=0.78in]{geometry}
\usepackage{array}
\usepackage{tabularx}
\usepackage{microtype}
\usepackage{xcolor}
\usepackage{hyperref}
\usepackage{fancyhdr}
\usepackage{lastpage}

\definecolor{navy}{HTML}{183756}
\definecolor{accent}{HTML}{2E74B5}
\definecolor{muted}{HTML}{5B656F}
\definecolor{rulegray}{HTML}{DADFE5}
\definecolor{ink}{HTML}{1C2126}

\hypersetup{
  colorlinks=true,
  urlcolor=accent,
  linkcolor=accent,
  pdfauthor={Shahil Shaik},
  pdftitle={Shahil Shaik --- Academic CV},
  pdfsubject={Academic CV for research computing allocation applications},
  pdfkeywords={robotics, reinforcement learning, transformers, VLA, multi-agent systems}
}
\urlstyle{same}

\pagestyle{fancy}
\fancyhf{}
\renewcommand{\headrulewidth}{0pt}
\fancyhead[R]{\color{muted}\fontsize{8}{9}\selectfont\bfseries SHAHIL SHAIK \enspace|\enspace ACADEMIC CV}
\fancyfoot[C]{\color{muted}\fontsize{8.5}{10}\selectfont Shahil Shaik \enspace$\bullet$\enspace Page \thepage\ of \pageref*{LastPage}}
\setlength{\headheight}{12pt}
\setlength{\parindent}{0pt}
\setlength{\parskip}{0pt}
\setlength{\emergencystretch}{3em}
\raggedbottom
\color{ink}

\newcommand{\cvsection}[1]{%
  \vspace{7pt}%
  {\color{navy}\fontsize{10.5}{12}\selectfont\bfseries\MakeUppercase{#1}}\par
  \vspace{2pt}\color{rulegray}\hrule height 0.55pt\color{ink}\vspace{4pt}
}

\newcommand{\researchentry}[3]{%
  {\color{navy}\bfseries #1}%
  {\color{muted}\enspace|\enspace\itshape #2}\par
  \vspace{1pt}#3\par\vspace{4pt}
}

\newcommand{\pubentry}[4]{%
  \hangindent=0.18in\hangafter=1
  #1 (#2). \textit{#3}. #4\par\vspace{5pt}
}

\newcommand{\skillrow}[2]{%
  {\color{navy}\bfseries #1} & #2\\[2.5pt]
}

\begin{document}
\raggedright

\begin{center}
  {\color{navy}\fontsize{22}{25}\selectfont\bfseries SHAHIL SHAIK}\par
  \vspace{4pt}
  {\color{muted}\fontsize{11.5}{14}\selectfont\bfseries
    Ph.D. Candidate in Mechanical Engineering \enspace|\enspace Robotics and Intelligent Systems}\par
  \vspace{5pt}
  {\fontsize{9.6}{11.5}\selectfont
    Clemson University
    \enspace$\bullet$\enspace
    \href{mailto:shahils@clemson.edu}{shahils@clemson.edu}
    \enspace$\bullet$\enspace
    \href{https://scholar.google.com/citations?user=THKUDwwAAAAJ&hl=en}{Google Scholar}
    \enspace$\bullet$\enspace
    \href{https://github.com/shahils01}{GitHub}}
\end{center}
\vspace{-1pt}

\cvsection{Research Profile}
Ph.D. candidate and Graduate Research Assistant developing efficient learning architectures
for embodied and multi-agent intelligence. Research spans transformer attention,
vision-language and vision-language-action models, multi-agent deep reinforcement learning,
distributional policy optimization, and shared human--robot control. Published in the
American Control Conference and \textit{IET Cyber-Systems and Robotics}; experienced in
reproducible, GPU-accelerated research with PyTorch.

\cvsection{Education}
\begin{tabularx}{\textwidth}{@{}>{\color{muted}\bfseries}p{1.25in}X@{}}
2021--Present &
  {\color{navy}\bfseries Ph.D. in Mechanical Engineering (Robotics)}\\[-1pt]
& Clemson University, Clemson, South Carolina\\[-1pt]
& {\color{muted}\small Expected graduation: August 2027. Research focus: intelligent multi-agent robotic systems and learning-based autonomy.}\\[5pt]
2019--2021 &
  {\color{navy}\bfseries Graduate Study in Mechanical Engineering}\\[-1pt]
& Clemson University, Clemson, South Carolina\\[-1pt]
& {\color{muted}\small Entered the M.S. program in 2019 and transitioned to the Ph.D. program in 2021.}\\[5pt]
2015--2019 &
  {\color{navy}\bfseries Bachelor's Degree in Mechanical Engineering}\\[-1pt]
& Anil Neerukonda Institute of Technology and Sciences, Visakhapatnam, India\\[-1pt]
& {\color{muted}\small GPA: 8.55/10.}
\end{tabularx}

\cvsection{Academic Appointment}
\begin{tabularx}{\textwidth}{@{}>{\color{muted}\bfseries}p{1.25in}X@{}}
Current &
  {\color{navy}\bfseries Graduate Research Assistant, Department of Mechanical Engineering}\\[-1pt]
& Clemson University, Clemson, South Carolina\\[-1pt]
& {\color{muted}\small Faculty advisor: Dr. Yue Wang.}
\end{tabularx}

\cvsection{Research Interests}
Efficient transformer architectures \enspace$\bullet$\enspace
Vision-language-action learning \enspace$\bullet$\enspace
Multi-agent reinforcement learning \enspace$\bullet$\enspace
Robot learning \enspace$\bullet$\enspace
Human--robot shared control \enspace$\bullet$\enspace
Distributed and memory-efficient deep learning

\cvsection{Selected Research}
\researchentry
  {Lie-Generated Metric Attention (LGMA)}
  {Lead developer and investigator}
  {Developing a PyTorch attention architecture that generates head-specific metrics from
  shared Lie-algebra generators to reduce parameter and key--value cache costs. Built
  controlled MHA-versus-LGMA experiments, stability diagnostics, checkpointing, efficiency
  accounting, and single-/multi-GPU training workflows for TinyStories and robot-learning studies.}

\researchentry
  {Vision-Language(-Action) Models for Multi-Agent Robotics}
  {Research and software development}
  {Investigating multimodal critics and action models for policy evaluation and robot learning,
  including MA-VLCM and planned LIBERO benchmarking. Work integrates visual observations,
  language-conditioned reasoning, temporal context, and multi-agent state representations.}

\researchentry
  {Multi-Agent Deep Reinforcement Learning}
  {Doctoral research}
  {Developed methods for learning under constrained inter-agent communication and studied
  distributional policy-gradient estimation. Research emphasizes coordination, robustness,
  sample efficiency, and principled evaluation across multi-agent settings.}

\researchentry
  {Shared Control of Mobile Robots}
  {Collaborative research}
  {Studied deep reinforcement-learning approaches that combine human guidance with autonomous
  control for mobile robot navigation, resulting in a peer-reviewed journal publication.}

\newpage

\cvsection{Publications and Preprints}
\pubentry
  {\textbf{\color{navy}Shaik, S.}, Smereka, J.~M., \& Wang, Y.}
  {2026}
  {Multi-Agent Deep Reinforcement Learning Under Constrained Communications}
  {arXiv preprint arXiv:2601.17069.}

\pubentry
  {\textbf{\color{navy}Shaik, S.}, Parameshwaran, A., Nayak, A., Smereka, J.~M., \& Wang, Y.}
  {2026}
  {MA-VLCM: A Vision Language Critic Model for Value Estimation of Policies in Multi-Agent Team Settings}
  {arXiv preprint arXiv:2603.15418.}

\pubentry
  {Nayak, A., \textbf{\color{navy}Shaik, S.}, \& Wang, Y.}
  {2026}
  {Belief-Aware VLM Model for Human-like Reasoning}
  {arXiv preprint arXiv:2604.09686.}

\pubentry
  {\textbf{\color{navy}Shaik, S.}, Smereka, J.~M., \& Wang, Y.}
  {2025}
  {Generalized Advantage Estimation for Distributional Policy Gradients}
  {Proceedings of the American Control Conference, 3073--3078.}

\pubentry
  {Tian, C., \textbf{\color{navy}Shaik, S.}, \& Wang, Y.}
  {2021}
  {Deep Reinforcement Learning for Shared Control of Mobile Robots}
  {\textit{IET Cyber-Systems and Robotics}, 3(4), 315--330.}

\cvsection{Research Software and Reproducibility}
{\color{navy}\bfseries Open-source LGMA implementation.}
Maintains a public PyTorch research prototype with interchangeable LGMA and conventional
multi-head attention, synthetic and TinyStories experiments, and VLA
data/model/evaluation modules.
\href{https://github.com/shahils01/Attention-is-Transformed}{Repository}\par\vspace{5pt}

{\color{navy}\bfseries Experiment integrity.}
Built unit and smoke tests, configuration-driven experiments, resumable checkpoints,
validation workflows, per-GPU throughput and memory accounting, and automated summaries
for validation-loss-versus-token and validation-loss-versus-GPU-hour comparisons.\par\vspace{5pt}

{\color{navy}\bfseries Research-computing readiness.}
Prepared bf16 CUDA training, gradient accumulation, single- and multi-GPU
DistributedDataParallel execution, Slurm launch patterns, dataset staging, and portable
Conda/Apptainer environment plans for accelerator-based studies.

\cvsection{Technical Skills}
\begin{tabularx}{\textwidth}{@{}p{1.48in}X@{}}
\skillrow{Programming}{Python; scientific computing and research software development}
\skillrow{Deep Learning}{PyTorch; transformer models; attention mechanisms; reinforcement learning; VLM/VLA}
\skillrow{GPU \& Scale}{CUDA; bf16 mixed precision; PyTorch DistributedDataParallel; torchrun; Slurm}
\skillrow{Research Stack}{Hugging Face Datasets/Transformers; NumPy; h5py; OpenCV; Weights \& Biases}
\skillrow{Robotics \& Tools}{MuJoCo/LIBERO workflows; Git/GitHub; Conda; Apptainer; pytest}
\end{tabularx}

\end{document}
